% Part: applied-modal-logics
% Chapter: temporal-logic
% Section: properties-accessibility

\documentclass[../../../include/open-logic-section]{subfiles}

\begin{document}

\olfileid[fr]{aml}{tl}{acc}

\olsection{Propriétés des cadres temporels}

Puisque nos modèles temporels n'imposent aucune condition à la
relation~$\prec$, le seul de nos axiomes usuels qui soit valide dans
tous les modèles est~$K$, ou ses analogues $K_{\Gtemp}$ et~$K_{\Htemp}$ :
\begin{align*}
\tag{$K_{\Gtemp}$}  \Gtemp (p \to q) & \to (\Gtemp p \to \Gtemp q)\\
\tag{$K_{\Htemp}$}  \Htemp (p \to q) & \to (\Htemp p \to \Htemp q)
\end{align*}

Si nous souhaitons imposer à nos modèles des conditions plus strictes
sur la représentation du temps, par exemple en exigeant que $\prec$
soit un ordre linéaire, nous obtiendrons d'autres formules valides
dans ces modèles.

\begin{table}[t]
  \begin{tabular}{| p{\dimexpr.5\linewidth-2\tabcolsep-2\arrayrulewidth-.5\doublerulesep\relax} || p{\dimexpr.5\linewidth-2\tabcolsep-2\arrayrulewidth-.5\doublerulesep\relax} |}
    \hline
    {\emph{Si $\prec$ est \dots}} & {\emph{alors \dots est vraie dans~$\mModel{M}$ :}} \\
    \hline \hline
    \emph{transitive} : \newline
    $\forall u \forall v \forall w ((u \prec v \land v \prec w) \lif u \prec w)$ &
    $\Ftemp \Ftemp p \lif \Ftemp p$  \\
    \hline
    \emph{connexe (comparabilité)} : \newline
    $\forall w \forall v (w \prec v \lor w = v \lor v \prec w)$ &
    $(\Ftemp \Ptemp  p \lor \Ptemp \Ftemp p) \lif (\Ptemp p \lor p \lor \Ftemp p)$\\
    \hline
    \emph{dense} : \newline
    $\forall w \forall v (w \prec v \to \exists u(w \prec u \land u \prec v))$ &
    $\Ftemp p \lif \Ftemp \Ftemp p$ \\
    \hline
    \emph{sans borne (vers le passé)} : \newline
    $\forall w \exists v( v \prec w)$ &
    $\Htemp p \to \Ptemp p$ \\
    \hline
    \emph{sans borne (vers le futur)} : \newline
    $\forall w \exists v( w \prec v)$ &
    $\Gtemp p \to \Ftemp p$ \\
    \hline
  \end{tabular}
  \caption{Quelques liens entre propriétés des cadres temporels et formules valides.}
  \ollabel{tab:correspondence}
\end{table}

Dans un cadre connexe, si $p$ est vrai à un état quelconque,
alors, à tout état, au moins l'une des formules $\Ptemp p$, $p$
et~$\Ftemp p$ est vraie. Cela explique l'implication à la deuxième
ligne du~\olref{tab:correspondence}. La réciproque échoue : dans
un cadre formé de deux états munis chacun d'une boucle et sans
flèche entre eux, $\Ptemp p$ et $\Ftemp p$ ont la même valeur de
vérité que~$p$ en chaque état. La formule du tableau y est donc
valide malgré l'absence de comparabilité entre les deux états.

L'irréflexivité $\forall w\,\lnot(w\prec w)$ ne se caractérise
par aucune formule de ce langage. Considérons le cycle à deux
états et le cadre à un seul état muni d'une boucle. Pour toute
valuation sur le second cadre, donnons à chaque variable la même
valeur aux deux états du cycle qu'à l'unique état du cadre à
boucle. Par induction sur les formules, chaque formule a alors
la même valeur aux trois états. Les clauses atomiques et
booléennes le donnent directement. Pour chacune des clauses
temporelles, chaque état a exactement un prédécesseur et un
successeur, et l'hypothèse d'induction s'applique à la sous-formule
en ces états. Toute formule valide sur le cycle est donc valide
sur le cadre à boucle. Or seul le cycle est irréflexif.
\end{document}
